% Montage d'un titrage avec suivi pH-métrique (L'essentiel, chapitre 3)
\begin{tikzpicture}[x=1cm,y=1cm,scale=0.92,every node/.style={transform shape=false}]
\pic at (0,0) {potence={hauteur=7.4}};
\pic at (0,0.16) {agitateur={largeur=2.3}};
\pic at (0,0.66) {becher={largeur=1.6,hauteur=1.9,niveau=1.15,liquide=pcEau}};
\pic at (-0.15,0.74) {barreau};
\pic at (0.02,2.75) {burette={hauteur=5.0,niveau=4.3,liquide=otNavy!14}};
\pic at (0,6.9) {pince};
\pic (sd) at (0.45,0.9) {sonde={hauteur=2.5}};
\pic (af) at (3.3,2.0) {afficheur={largeur=1.8,hauteur=1.2,texte={7,52}}};
\draw[pcFil] (sd-haut) to[out=90,in=180] (af-prise);
\draw[pcRepere] (0.14,6.0) -- (1.2,6.4) node[pcLegende,anchor=west]{burette graduée :\\solution titrante ($c_\mathrm{B}$ connue)};
\draw[pcRepere] (0.45,3.15) -- (1.3,4.05) node[pcLegende,anchor=west]{robinet};
\draw[pcRepere] (-0.55,1.25) -- (-1.75,1.75) node[pcLegende,anchor=east]{solution titrée :\\volume $V_\mathrm{A}$ prélevé\\à la pipette jaugée};
\draw[pcRepere] (0.5,2.6) -- (1.3,3.62) node[pcLegende,anchor=west]{sonde du pH-mètre};
\draw[pcRepere] (-1.0,0.4) -- (-1.75,0.55) node[pcLegende,anchor=east]{agitateur magnétique};
\node[pcLegende,anchor=north] at (3.3,1.95) {pH-mètre étalonné};
\end{tikzpicture}
\par\smallskip\legende{Pour un suivi conductimétrique, la sonde est remplacée par une cellule de conductimétrie reliée à un conductimètre.}
